% ============================================================
% MẪU BÁO CÁO THỰC TẬP KỸ THUẬT / ĐỒ ÁN TỐT NGHIỆP
% ĐẠI HỌC BÁCH KHOA HÀ NỘI - KHOA TOÁN TIN
% Phụ lục I: Mẫu bìa báo cáo
% Biên dịch bằng XeLaTeX
% ============================================================

\documentclass[12pt,a4paper]{report}

% ========== PACKAGES ==========
\usepackage{fontspec}
\defaultfontfeatures{Ligatures=TeX}
\setmainfont{Times New Roman}
\setsansfont{Arial}
\setmonofont{Courier New}
\usepackage[vietnamese]{babel}
\usepackage{geometry}
\usepackage{graphicx}
\usepackage{hyperref}
\usepackage{indentfirst}
\usepackage{setspace}
\usepackage{float}
\usepackage{caption}
\usepackage{enumitem}
\usepackage{makeidx}
\usepackage{listings}
\usepackage{xcolor}
\usepackage{booktabs}
\usepackage{amsmath,amssymb}
\usepackage{longtable}

% ========== HYPERREF ==========
\hypersetup{
    colorlinks=true,
    linkcolor=black,
    citecolor=black,
    urlcolor=black,
    pdfborder={0 0 0}
}

% ========== CĂN LỀ (A4, khổ chuẩn) ==========
\geometry{
    left=2.5cm,
    right=2cm,
    top=2.5cm,
    bottom=2.5cm
}

% ========== GIÃN DÒNG ==========
\onehalfspacing

% ========== THỤT ĐẦU DÒNG ==========
\setlength{\parindent}{1.27cm}
\setlength{\parskip}{0pt}

% ========== CHỈ MỤC ==========
\makeindex

% ========== CODE LISTING ==========
\lstset{
    basicstyle=\ttfamily\small,
    breaklines=true,
    frame=single,
    numbers=left,
    numberstyle=\tiny,
    commentstyle=\color{gray},
    keywordstyle=\color{blue},
    stringstyle=\color{red},
    showstringspaces=false,
    tabsize=2
}

% ========== THÔNG TIN ĐỒ ÁN - ĐIỀN VÀO ĐÂY ==========
\newcommand{\tenDeTai}{Phân tích thiết kế hệ thống chấm điểm tự động}
\newcommand{\tenSinhVien}{NGUYỄN VĂN A}
\newcommand{\emailSinhVien}{nguyenvanabc@sis.hust.edu.vn}
\newcommand{\maSinhVien}{2024666888}
\newcommand{\chuyenNganh}{Toán -- Tin / Hệ thống Thông tin Quản lý}
\newcommand{\tenGVHD}{PGS. TS. Phạm Văn ABC}
\newcommand{\thangNam}{6/2024}
\newcommand{\loaiBaoCao}{THỰC TẬP KỸ THUẬT / ĐỒ ÁN TỐT NGHIỆP (1,2..)}

% ============================================================
\begin{document}
\fontsize{13}{16}\selectfont  % Cỡ chữ mặc định 13pt như PDF gốc
\pagestyle{empty}  % Không đánh số trang

% ============================================================
% TRANG 1: BÌA BÁO CÁO
% ============================================================
\begin{center}

    {\fontsize{13}{16}\selectfont Phụ lục I: Mẫu bìa báo cáo}

    \vspace{0.5cm}

    {\fontsize{14}{18}\selectfont\bfseries ĐẠI HỌC BÁCH KHOA HÀ NỘI}\\[4pt]
    {\fontsize{14}{18}\selectfont\bfseries KHOA TOÁN -- TIN}

    \vspace{5cm}

    {\fontsize{18}{22}\selectfont\bfseries \loaiBaoCao}

    \vspace{1.5cm}

    {\fontsize{14}{18}\selectfont\bfseries \tenDeTai}

    \vspace{1cm}

    {\fontsize{14}{18}\selectfont\bfseries \tenSinhVien}

    \vspace{0.6cm}

    {\fontsize{14}{18}\selectfont Email: \emailSinhVien}\\[4pt]
    {\fontsize{14}{18}\selectfont Mã sinh viên: \maSinhVien}

    \vspace{0.6cm}

    {\fontsize{14}{18}\selectfont\bfseries Chuyên ngành \chuyenNganh}

    \vfill

    \begin{flushleft}
        \hspace{2cm}
        \begin{minipage}{0.6\textwidth}
            {\fontsize{13}{16}\selectfont\bfseries Giảng viên hướng dẫn: \tenGVHD}\\[8pt]
            \hspace{5cm}{\fontsize{10}{12}\selectfont Chữ ký của GVHD}
        \end{minipage}
    \end{flushleft}

    \vspace{1.5cm}

    {\fontsize{13}{16}\selectfont\bfseries HÀ NỘI, \thangNam}

\end{center}

\newpage

% ============================================================
% TRANG 2: BÌA PHỤ
% ============================================================
\begin{center}

    \vspace*{1cm}

    {\fontsize{14}{18}\selectfont\bfseries ĐẠI HỌC BÁCH KHOA HÀ NỘI}\\[4pt]
    {\fontsize{14}{18}\selectfont\bfseries KHOA TOÁN -- TIN}

    \vspace{5cm}

    {\fontsize{18}{22}\selectfont\bfseries \loaiBaoCao}

    \vspace{1.5cm}

    {\fontsize{14}{18}\selectfont\bfseries \tenDeTai}

    \vspace{1cm}

    {\fontsize{14}{18}\selectfont\bfseries \tenSinhVien}

    \vspace{0.6cm}

    {\fontsize{14}{18}\selectfont Email: \emailSinhVien}\\[4pt]
    {\fontsize{14}{18}\selectfont Mã sinh viên: \maSinhVien}

    \vspace{0.6cm}

    {\fontsize{14}{18}\selectfont\bfseries Chuyên ngành \chuyenNganh}

    \vfill

    \begin{flushleft}
        \hspace{2cm}
        \begin{minipage}{0.6\textwidth}
            {\fontsize{13}{16}\selectfont\bfseries Giảng viên hướng dẫn: \tenGVHD}\\[8pt]
            \hspace{5cm}{\fontsize{10}{12}\selectfont Chữ ký của GVHD}
        \end{minipage}
    \end{flushleft}

    \vspace{1.5cm}

    {\fontsize{13}{16}\selectfont\bfseries HÀ NỘI, \thangNam}

\end{center}

\newpage

% ============================================================
% TRANG 3: NHẬN XÉT CỦA GIẢNG VIÊN HƯỚNG DẪN
% ============================================================
\begin{center}
    {\fontsize{13}{16}\selectfont\bfseries NHẬN XÉT CỦA GIẢNG VIÊN HƯỚNG DẪN}
\end{center}

\vspace{0.8cm}

Biểu mẫu của Đề tài/khóa luận tốt nghiệp theo qui định của viện, tuy nhiên cần đảm bảo giáo viên giao đề tài ký và ghi rõ họ và tên.

Trường hợp có 2 giáo viên hướng dẫn thì sẽ cùng ký tên.

\vspace{9cm}

\begin{flushright}
    \begin{minipage}{0.35\textwidth}
        \centering
        Hà Nội, ngày \dots.\ tháng \dots\ \ năm \dots\\[12pt]
        {\fontsize{13}{16}\selectfont Giáo viên hướng dẫn}\\[20pt]
        {\fontsize{10}{12}\selectfont Ký và ghi rõ họ tên}
    \end{minipage}
\end{flushright}

\newpage

% ============================================================
% TRANG 4: PHIẾU BÁO CÁO TIẾN ĐỘ ĐỒ ÁN
% ============================================================
\begin{center}
    {\fontsize{13}{16}\selectfont\bfseries PHIẾU BÁO CÁO TIẾN ĐỘ ĐỒ ÁN}
\end{center}

\vspace{1.5cm}

\begin{center}
    {\fontsize{13}{16}\selectfont - Phiếu này in từ hệ thống số hóa}
\end{center}

\newpage

% ============================================================
% TRANG 5: LỜI CẢM ƠN + TÓM TẮT NỘI DUNG ĐỒ ÁN
% ============================================================
\begin{center}
    {\fontsize{13}{16}\selectfont\bfseries Lời cảm ơn}
\end{center}

\vspace{0.8cm}

Đây là mục tùy chọn, nên viết phần cảm ơn ngắn gọn, tránh dùng các từ sáo rỗng, giới hạn trong khoảng 100-150 từ.

% (Sinh viên tự viết nội dung cảm ơn tại đây)

\vspace{1.5cm}

\begin{center}
    {\fontsize{13}{16}\selectfont\bfseries Tóm tắt nội dung đồ án}
\end{center}

\vspace{0.8cm}

Tóm tắt nội dung của đồ án tốt nghiệp trong khoảng tối đa 300 chữ. Phần tóm tắt cần nêu được các ý: vấn đề cần thực hiện; phương pháp thực hiện; công cụ sử dụng (phần mềm, phần cứng\dots); kết quả của đồ án có phù hợp với các vấn đề đã đặt ra hay không; tính thực tế của đồ án, định hướng phát triển mở rộng của đồ án (nếu có); các kiến thức và kỹ năng mà sinh viên đã đạt được.

% (Sinh viên tự viết nội dung tóm tắt tại đây)

\vspace{3cm}

\begin{flushright}
    \begin{minipage}{0.35\textwidth}
        \centering
        Hà Nội, ngày \dots.\ tháng \dots\ năm\dots\\[12pt]
        {\fontsize{13}{16}\selectfont Sinh viên thực hiện}\\[20pt]
        {\fontsize{10}{12}\selectfont Ký và ghi rõ họ tên}
    \end{minipage}
\end{flushright}

\newpage

% ============================================================
% TRANG 6: MỤC LỤC
% ============================================================
\begin{center}
    {\fontsize{13}{16}\selectfont\bfseries MỤC LỤC}
\end{center}

\vspace{0.8cm}

% (Mục lục sẽ được sinh viên tự điền hoặc dùng lệnh \tableofcontents)

\newpage

% ============================================================
% TRANG 7: DANH MỤC HÌNH VẼ
% ============================================================
\begin{center}
    {\fontsize{13}{16}\selectfont\bfseries DANH MỤC HÌNH VẼ}
\end{center}

\vspace{0.8cm}

% (Danh sách hình vẽ - sinh viên tự điền)

\newpage

% ============================================================
% TRANG 8: DANH MỤC KÝ HIỆU
% ============================================================
\begin{center}
    {\fontsize{13}{16}\selectfont\bfseries DANH MỤC KÝ HIỆU}
\end{center}

\vspace{0.8cm}

% (Danh sách ký hiệu - sinh viên tự điền)

\newpage

% ============================================================
% TRANG 9: CHƯƠNG 1, 2, 3 (MẪU CẤU TRÚC CHƯƠNG)
% ============================================================
{\fontsize{13}{16}\selectfont\bfseries CHƯƠNG 1. GIỚI THIỆU}

\vspace{0.8cm}

% Nội dung chương 1...

\vspace{0.3cm}

{\fontsize{13}{16}\selectfont\bfseries CHƯƠNG 2.}

\vspace{0.8cm}

% Nội dung chương 2...

\vspace{0.3cm}

{\fontsize{13}{16}\selectfont\bfseries CHƯƠNG 3.}

\vspace{0.8cm}

% Nội dung chương 3...

\newpage

% ============================================================
% TRANG 10: KẾT LUẬN
% ============================================================
\begin{center}
    {\fontsize{13}{16}\selectfont\bfseries KẾT LUẬN}
\end{center}

\vspace{0.8cm}

% (Sinh viên tự viết kết luận)

\newpage

% ============================================================
% TRANG 11: PHỤ LỤC
% ============================================================
\begin{center}
    {\fontsize{13}{16}\selectfont\bfseries PHỤ LỤC}
\end{center}

\vspace{0.8cm}

Phần này bao gồm những nội dung cần thiết nhằm minh họa hoặc hỗ trợ cho nội dung đồ án như số liệu, mã nguồn chương trình, biểu mẫu, tranh ảnh, thông số kỹ thuật, \dots. Phụ lục không được dày hơn phần chính của đồ án.

Các nội dung khác nhau được đặt trên các Phụ Lục khác nhau kèm tựa tương ứng. Thứ tự Phụ lục đánh tên theo A, B, C, \dots

Ví dụ:

\hspace{1cm}{\fontsize{13}{16}\selectfont Phụ lục A \quad MÃ NGUỒN CHƯƠNG TRÌNH}

\hspace{1cm}{\fontsize{13}{16}\selectfont Phụ lục B \quad HƯỚNG DẪN SỬ DỤNG SẢN PHẨM THI CÔNG}

\newpage

% ============================================================
% TRANG 12: CHỈ MỤC
% ============================================================
\begin{center}
    {\fontsize{13}{16}\selectfont\bfseries CHỈ MỤC}
\end{center}

\vspace{0.8cm}

% (Chỉ mục - sinh viên tự điền hoặc dùng \printindex)

\newpage

% ============================================================
% TRANG 13: TÀI LIỆU THAM KHẢO
% ============================================================
\begin{center}
    {\fontsize{13}{16}\selectfont\bfseries TÀI LIỆU THAM KHẢO}
\end{center}

\vspace{0.8cm}

% (Sinh viên tự điền danh sách tài liệu tham khảo)
% Ví dụ:
% [1] Tác giả A, \textit{Tên sách}, Nhà xuất bản, Năm.
% [2] Tác giả B, \textit{Tên bài báo}, Tạp chí, Năm.

\end{document}
